\section{Reflected forms and compact tests}
\label{sec:basics}

Throughout, a divisor is a locally finite multiset of complex numbers.
We sum over its distinct locations and denote their positive integer
multiplicities by \(\nu(z)\). Fix a strip half-width \(a>0\), and
assume
\begin{equation}\label{eq:divisor}
 |\Re z|\le a,\qquad z^\#=-\overline z,\qquad
 z\in Z\Longleftrightarrow z^\#\in Z,\qquad
 \nu(z)=\nu(z^\#).
\end{equation}
Unless explicitly restricted to a finite divisor, we assume a polynomial
ordinate count: there are finite \(C>0,m\ge0\) such that
\begin{equation}\label{eq:polynomial-count}
 N_0(Y):=\sum_{|\Im z|\le Y}\nu(z)\le C(1+Y)^m
 \qquad(Y\ge0).
\end{equation}
The constants need not be uniform over all divisors except in the
uniform assertions where this is stated.

For \(f,g\in C_c^\infty(\R;\C)\), set
\begin{equation}\label{eq:form}
 H_f(z)=\int_\R f(u)e^{zu}\,du,\qquad
 Q_Z(f,g)=\sum_{z\in Z}\nu(z)H_f(z)\overline{H_g(z^\#)}.
\end{equation}
Our inner products are linear in the first variable. Repeated
integration by parts gives, for every integer \(j\ge0\),
\[
 \sup_{|\sigma|\le a}|H_f(\sigma+it)|
 \le C_{f,j,a}(1+|t|)^{-j}.
\]
For \(|t|\ge1\), this follows by integrating the Fourier transform
of \(e^{\sigma u}f(u)\) by parts \(j\) times, uniformly in the
bounded real parameter; bounded \(t\) is immediate. Dyadic summation
using \eqref{eq:polynomial-count} proves absolute convergence of
\eqref{eq:form}. Changing \(z\) to \(z^\#\) then proves
\(Q_Z(f,g)=\overline{Q_Z(g,f)}\).

For \(z=\sigma+i\gamma\) with \(\sigma>0\), write
\[
 A_z(f)=\frac{H_f(\sigma+i\gamma)+H_f(-\sigma+i\gamma)}2,
 \qquad
 D_z(f)=\frac{H_f(\sigma+i\gamma)-H_f(-\sigma+i\gamma)}2.
\]
The two locations of this pair contribute exactly
\begin{equation}\label{eq:pair}
 2\nu(z)\bigl(|A_z(f)|^2-|D_z(f)|^2\bigr).
\end{equation}
An axis location \(i\gamma\) contributes
\(\nu(i\gamma)|H_f(i\gamma)|^2\ge0\). Formula
\eqref{eq:pair} explains both the possibility of a negative test
and the possibility that positive terms screen it. Multiplicity
weights a value; it does not ask the transform to vanish to higher
order.

The quantities studied below are
\begin{equation}\label{eq:envelope-radius}
 \begin{split}
 \cM_Z(L)&=\sup_{0\ne f\in C_c^\infty(-L,L)}
       \left(-Q_Z(f,f)/\norm f_2^2\right)_+,\qquad L>0,\\
 r_Z(\eta)&=\inf\{L>0:\ \exists\,0\ne f\in C_c^\infty(-L,L),
                 \ Q_Z(f,f)\le-\eta\norm f_2^2\},\qquad\eta>0.
 \end{split}
\end{equation}
Here \(x_+=\max(x,0)\), the first supremum may be infinite,
and the infimum of the empty set is infinite. No extremum is assumed
attained. In particular, \(\cM_Z(L)>\eta\) supplies an actual
test with margin exceeding \(\eta\), whereas equality alone need
not. Enlarging \(L\) enlarges the test class, so \(\cM_Z\)
is nondecreasing. Adding an axis divisor can only decrease this
envelope. We say that the form is locally semibounded when
\(\cM_Z(L)<\infty\) for every finite \(L>0\); this terminology
asserts a bound on the smooth test class, not a closed operator
realization.

\subsection{Changing the center and the scale}
For a target ordinate \(\gamma\), use the centered count
\[
 N_\gamma(Y)=\sum_{|\Im z-\gamma|\le Y}\nu(z).
\]
If \(W=Z-i\gamma\) and \(f(u)=e^{-i\gamma u}v(u)\), then
\(H_f(z)=H_v(z-i\gamma)\) and
\(Q_Z(f,f)=Q_W(v,v)\). Norm and support are unchanged.
Thus a target can be moved to the real axis without a
height-dependent cost in the test norm. An origin-centered count
\eqref{eq:polynomial-count} implies the centered bound
\[
 N_\gamma(Y)\le C(1+|\gamma|)^m(1+Y)^m.
\]
The constant here depends on the chosen target.

For \(\alpha>0\), let \(W=Z/\alpha\), with the same
multiplicities, and put \(f(u)=\sqrt\alpha\,v(\alpha u)\).
Direct substitution yields
\begin{equation}\label{eq:scale}
 \norm f_2=\norm v_2,\qquad
 H_f(z)=\alpha^{-1/2}H_v(z/\alpha),\qquad
 Q_Z(f,f)=\alpha^{-1}Q_W(v,v).
\end{equation}
A radius \(L\) for \(v\) becomes \(L/\alpha\) for \(f\).
These factors will matter when passing from a normalized strip to
a growth rate.

\subsection{Finite features and their inertia}
For a finite divisor, \eqref{eq:form} extends continuously to
\(L^2(-L,L)\). List its distinct locations as \(z_1,\ldots,z_N\),
and define \((Ef)_j=\sqrt{\nu(z_j)}H_f(z_j)\).
Let \(J\) permute coordinates according to reflection. Then
\begin{equation}\label{eq:finite-features}
 Q_Z(f,f)=\langle Ef,JEf\rangle_{\C^N},\qquad
 A_Z=E^*JE.
\end{equation}
The exponential functions representing the evaluations are linearly
independent on every nonempty interval. Indeed, a vanishing finite
linear combination is analytic and hence vanishes identically;
its first \(N\) derivatives at any point give an invertible
Vandermonde system. Thus \(E\) is onto. On the orthogonal complement
of its kernel, \eqref{eq:finite-features} is a congruence of \(J\).
There is one negative direction for each distinct off-axis reflected
pair, one positive direction for each such pair, and one positive
direction for each distinct axis location. All remaining directions
are null. This finite-inertia mechanism is classical; compare
Bombieri's finite exponential-kernel analysis \cite[Sections 8--9]{Bombieri}.

Smooth compact functions are dense in \(L^2(-L,L)\), and the
finite form is continuous. Consequently its least Rayleigh value
and negative envelope can be computed on either class. A strict
inequality found in \(L^2\) survives smooth approximation. For an
infinite axis background the form need not be bounded, so this
density argument cannot be used without additional estimates.
The fixed-background argument below supplies those estimates.

Finally, for the single unit pair \([\beta]+[-\beta]\),
\eqref{eq:pair} and the even/odd decomposition give
\begin{equation}\label{eq:single-pair}
 \cM_{[\beta]+[-\beta]}(L)
 =2\norm{\sinh(\beta u)}_{L^2(-L,L)}^2
 =\beta^{-1}\{\sinh(2\beta L)-2\beta L\}.
\end{equation}
This exact elementary example has size
\((4/3)\beta^2L^3\) when \(\beta L\) tends to zero and
size \(e^{2\beta L}/(2\beta)\) when \(\beta L\) tends to
infinity. The question is how much these scales can be delayed
by a counting budget and a positive background.
